%% 07_conclusion.tex
\section{Conclusion}
\label{sec:conclusion}

This paper documents a browser-based replication of the
experimental asset market of~\citet{dufwenberg2005bubbles},
extended with the expected-utility messaging population
of~\citet{lopezlira2025llm}, and introduces a new paradigm in
which a single ChatGPT call seeds each utility agent's private
subjective valuation at run~start. The innovation is deliberately
minimal: one number per agent, clamped to
$[0.25\,\FV_0,\,1.75\,\FV_0]$, written to the
\texttt{psychAnchor} slot and consumed on the first decision tick
of the standard Lopez-Lira prior. Everything downstream of the
first tick---the belief blend, the trust EMA, the expected-utility
action choice, the order matching, the dividend crediting, and
the market-quality instruments---is unchanged from the replicated
baseline, and the simulator's seeded mulberry32 generator ensures
that every run can be replayed identically with or without network
access.

The contribution is threefold. First, the testbed constitutes a
clean vanilla-JavaScript implementation of two otherwise unrelated
experimental designs, specified in enough detail
(Sections~\ref{sec:design}--\ref{sec:agents}) that a reader can
reconstruct the model from the paper alone. Second, the Wang~2026
innovation introduces a narrowly-scoped use of a language model as
a \emph{psychological prior elicitor} rather than as a trader: the
model is never asked to choose an action, and its output is routed
through the same prior slot used by the deterministic fallback,
so that any behavioural effect is attributable to the starting
valuation rather than to an expansion of the action space. Third,
the research question---does the asset end up in the hands of the
agent with the highest anchored valuation---is instrumented with
a per-tick match indicator and a normalised gap, and reported
alongside the classical bubble statistics
(Haessel~$R^2$, normalised deviation, amplitude, turnover,
allocative efficiency). The instrumentation is designed so that
any replication of the Wang~2026 paradigm can report the same
panel without additional~code.

Three caveats shape the preliminary findings. First, the ranges in
Table~\ref{tab:mktq} are illustrative, not calibrated: a full
empirical test requires many replications per cell, a pooled
treatment of the network variance of the anchor, and a formal
baseline against a deterministic prior with matched mean and
variance. Second, the anchor is consumed exactly once, which bounds
the downstream signal at whatever the first-tick prior contributes
to the five-action expected-utility score; refreshing $\hat V_i$
each period would be a natural extension but is in tension with the
reproducibility invariant. Third, generalising beyond the DLM~2005
dividend distribution and horizon would require re-deriving the
market-quality instruments and re-calibrating the experienced
profiles of Table~\ref{tab:exp}.

The testbed is free of build steps and of third-party dependencies,
and ships with the architecture diagram, glossary, and slide deck
used in this paper. Its ``Experiment settings'' panel exposes every
numerical constant that shapes a run---the population mix, the
risk-preference share, the passive fill probability, the trust
learning rate, the bias amplitude, and the valuation noise---and
the seeded PRNG makes every edit to a constant observable as a
difference against a matched baseline. The Wang~2026 paradigm is
the first extension of this infrastructure that asks a question
about language models, and it opens several follow-ups: whether
the anchor's effect on allocative efficiency scales with the
cross-agent variance of~$\hat V$; whether multiple language models
produce systematically different anchor vectors for the same
prompt; and whether adding deceptive messaging amplifies or dampens
the first-tick stratification that the anchor produces.
